\section{First-pass method}\label{sec:method}

This section describes the first pass, which runs the declared first
stages against public aggregates for all five episodes; the energy
episode is additionally run on household records in
Section~\ref{sec:results-secondstage}, whose method is set out there. Everything is generated by a pinned
pipeline that verifies its raw inputs by hash and emits every reported
value as a machine-generated macro or generated table; the declared
first-stage constants of Section~\ref{sec:episodes} are typed
\emph{with their citations} and marked as imported, and a small number
of external anchors and convention settings (cap levels, scheme
amounts, sensitivity dials) appear as literals beside the macros they
condition.

Following the companion studies' convention, the epistemic status of
every input is stated once:

\begin{center}
\small
\begin{tabular}{P{0.30\textwidth}P{0.60\textwidth}}
\toprule
Object & Status \\
\midrule
Episode first stages & \textbf{Imported} from cited published estimates (episodes 1--2) or government projections (episodes 4--5); episode 3 from the companion study. Never re-estimated here. \\
Expenditure by category $\times$ decile & \textbf{Data} (ONS Family Spending detailed workbooks). \\
Energy price-cap levels, statutory benefit rates, CPI outturns & \textbf{Data} (Ofgem/DESNZ, DWP, ONS). \\
Pass-through and utilisation dials & \textbf{Assumed}, bracketed, entering the pipeline once each. \\
Incidence arithmetic & \textbf{Computed}; no estimation step anywhere in the pipeline. \\
\bottomrule
\end{tabular}
\end{center}

\paragraph{Consumption incidence.} Each episode's price vector is
applied to the relevant categories of household expenditure by
equivalised-disposable-income decile, giving a per-decile burden or
gain in pounds a year, as a share of total spending, and as a share of
disposable income. These are first-order calculations: no substitution,
no behavioural response, fixed baskets, upper bounds on avoidable
welfare cost, in the envelope-theorem sense of \citet{deaton1989},
whose first-order incidence method this literature, from
\citet{fajgelbaumkhandelwal2016} onward, descends from.
Where a published aggregate exists, its role differs by episode and is
stated in the results: the OBR's Energy Price Guarantee cost is a
context anchor beside the pipeline's own aggregate, while the TCA's
\pounds 250 published cumulative is used as an \emph{anchor}: the
pipeline's linear-ramp cumulative overshoots it by a factor of
\MsTcaCrossCheckRatio,
so the headline TCA path is rescaled to the published figure
(Section~\ref{sec:results}), so that episode's level is the published
estimate's, not the pipeline's.

\paragraph{Rulebook arithmetic.} The vintage axis is seeded with
statutory arithmetic on public parameters: the real-terms shortfall of
the Universal Credit standard allowance during the 2022--23 inflation
episode relative to contemporaneous indexation (the uprating-lag cost,
computed month by month against CPI outturns), and the discrete
removal of the \pounds 20-a-week uplift in October 2021, the rulebook
event that immediately preceded the terms-of-trade shock. The
Energy Price Guarantee enters as the discretionary layer of episode 2:
the gap between the gross price shock and the EPG-truncated shock is
the discretionary cushioning share, computable per decile from cap
levels alone.

\paragraph{Two data caveats.} The FYE2025 Family Spending workbook
carries an ONS correction notice affecting percentage standard errors
across the high-level COICOP categories; the point estimates this paper
uses are unaffected, but no sampling-uncertainty statement is made on
the expenditure inputs for that reason. And the two vintages used
(FYE2022 for the TCA and energy episodes, whose spending bases must
predate the shocks, and FYE2025 for the India agreement) sit either
side of an ONS change in the equivalence scale, so income levels are
not comparable across vintages, which is why every cross-episode
comparison in this paper runs on shares of spending rather than shares
of income.

\paragraph{Comparability.} Episodes are compared in ratios, not
levels: burden as a share of spending by decile, cushioning shares
where computable, and a per-\pounds-billion-of-shock normalisation,
because gross sizes span more than two orders of magnitude and pound
totals would reduce the comparison to a statement about the energy
episode. Scenario dials (pass-through, utilisation) are reported as
labelled variants, never blended into central cells; brackets are outer
bounds, not confidence intervals; and estimate-based versus
projection-based first stages are flagged in every table.
